\chapter{Messagerie et architecture pilotée par les événements}\label{ch:pl:messaging-and-event-driven-architecture}

Un consommateur du topic des exécutions plante après avoir appliqué un lot de messages aux positions et avant d'enregistrer jusqu'où il a lu. Au
redémarrage, il relit le lot et l'applique de nouveau, si bien que la position de la société sur une action est doublée jusqu'à ce que quelqu'un la compare au
drop copy. Rien n'était cassé~: le système de messagerie a livré chaque message au moins une fois, ce qu'il avait exactement promis. Ce chapitre
porte sur ce que promettent les systèmes de messagerie, sur ce qu'ils ne peuvent pas promettre, et sur la manière dont un consommateur transforme une \omterm{def:pl:messaging-and-event-driven-architecture:delivery}{livraison au moins une fois} en effets exactement une fois --- avec un
\omterm{def:pl:messaging-and-event-driven-architecture:log}{journal partitionné} construit en mémoire, un \omterm{def:pl:position-and-pnl-services:service}{service de positions} soumis à mille plantages dans trois modes, et la \omterm{def:pl:messaging-and-event-driven-architecture:outbox}{boîte d'envoi transactionnelle} (outbox) qui permet à un service de modifier
sa base de données et de l'annoncer sans perdre ni l'un ni l'autre.

\section{Pourquoi les systèmes communiquent par messages}

La carte de la plateforme du chapitre~\ref{ch:pl:the-platform-map} est un ensemble de services~: \omterm{def:pl:trade-capture-and-booking:capture}{saisie des transactions} (trade capture), positions, risque, P\&L, surveillance. Ils peuvent
s'appeler directement, chacun sachant qui a besoin de quoi~; ou bien chacun peut publier ce qui s'est passé, et quiconque en a besoin s'y abonne. La seconde
option les découple de trois façons --- dans la connaissance (l'éditeur ne connaît pas ses consommateurs), dans le temps (un consommateur peut être arrêté et rattraper son retard), et
dans le débit (un consommateur lent ne ralentit pas l'éditeur) --- et fait du flux d'événements lui-même un enregistrement qu'on peut rejouer.

\begin{definition}[Publication--abonnement]\label{def:pl:messaging-and-event-driven-architecture:pubsub}
Le \emph{publication--abonnement}\index{publication--abonnement} (publish--subscribe) est un schéma de messagerie dans lequel des producteurs publient des messages dans des topics nommés sans
adresser de consommateur, et chaque consommateur reçoit les messages des topics auxquels il est abonné.
\end{definition}

Le coût est que le graphe des appels devient invisible, et que les pannes passent de l'appel (qui échoue bruyamment) à la livraison (qui peut perdre ou
répéter des messages en silence). La plus grande partie de ce chapitre traite de la seconde.

\section{Courtiers et transports sans courtier}

\begin{definition}[Courtier de messages, messagerie sans courtier]\label{def:pl:messaging-and-event-driven-architecture:broker}
Un \emph{courtier de messages}\index{courtier de messages} (message broker) est un serveur qui reçoit des messages de producteurs, les stocke et les livre à des consommateurs,
de sorte que producteurs et consommateurs ne connaissent que le courtier. Dans la \emph{messagerie sans courtier}\index{messagerie sans courtier} (brokerless messaging), les producteurs envoient directement
aux consommateurs --- en unicast, en multicast ou par mémoire partagée --- au moyen d'une bibliothèque présente dans chaque processus, sans serveur sur le chemin.
\end{definition}

Les deux répondent à des besoins différents. Un courtier ajoute un saut et un serveur à exploiter, et en contrepartie offre le stockage, le rejeu et des consommateurs qui vont et viennent.
Les transports sans courtier offrent la latence la plus faible et aucun composant central~: les chemins des données de marché et des ordres du livre~13 sont sans courtier, avec
du multicast (livre~13, chapitre~16) et des tampons circulaires (ring buffers) en mémoire partagée, et des bibliothèques comme Aeron --- «~efficient reliable UDP
unicast, UDP multicast, and IPC message transport~» --- assurent la fiabilité par-dessus. Entre services qui tiennent davantage à
ne pas perdre un événement qu'à la microseconde --- \omterm{def:pl:trade-capture-and-booking:capture}{saisie des transactions}, positions, risque, la chaîne post-marché --- la société utilise un journal.

\section{Le journal comme colonne vertébrale}

\begin{definition}[Journal partitionné, position du consommateur]\label{def:pl:messaging-and-event-driven-architecture:log}
Un \emph{journal partitionné}\index{journal partitionné} (partitioned log) stocke chaque topic sous forme de plusieurs partitions, chacune étant une suite de messages en ajout seul, repérés par
leur position dans celle-ci~; la partition d'un message est choisie à partir de sa clé, de sorte que des messages de même clé conservent leur ordre. Une
\emph{position du consommateur}\index{position du consommateur} (offset) est la position jusqu'à laquelle un consommateur a lu une partition, enregistrée par le consommateur (validée) afin de
pouvoir reprendre à cet endroit.
\end{definition}

\begin{definition}[Groupe de consommateurs]\label{def:pl:messaging-and-event-driven-architecture:group}
Un \emph{groupe de consommateurs}\index{groupe de consommateurs} (consumer group) est un ensemble de processus consommateurs qui se partagent les partitions d'un topic, chaque partition étant lue par exactement
un seul membre à la fois, avec une seule position validée par partition pour l'ensemble du groupe~; lorsque des membres rejoignent ou quittent le groupe, les partitions sont réaffectées
(un rééquilibrage).
\end{definition}

Telle est la conception que Kreps, Narkhede et Rao ont décrite pour Kafka en 2011~: un message n'a pas d'identifiant propre, «~each message is
addressed by its logical offset in the log~». Le \texttt{firm.eventlog} du chapitre la construit en mémoire~: des topics, des partitions selon un
hachage stable de la clé, des segments en ajout seul sur disque, une rétention qui supprime des segments entiers, et des positions validées par groupe. En regroupant les
exécutions par compte, toutes les exécutions d'un compte se retrouvent dans une même partition, dans l'ordre~; huit partitions partagées par trois consommateurs sont réparties en trois, trois et
deux.

\begin{omfigure}
\begin{tikzpicture}[font=\small,
  cell/.style={draw=omDef,fill=omDef!8,minimum width=6mm,minimum height=6mm,inner sep=0pt,font=\scriptsize},
  arr/.style={-{Stealth[length=2mm]},thick}]
\node[draw=omDef,rounded corners=2pt,align=center,text width=17mm,minimum height=24mm,font=\footnotesize] (prod) at (-0.9,0) {producteurs\\(saisie des\\trans\-actions)};
\foreach \p/\n in {0/7,1/5,2/8,3/6} {
  \node[font=\scriptsize,anchor=east] at (2.4,1.35-0.9*\p) {partition \p};
  \foreach \i in {1,...,\n} \node[cell] at (2.3+0.62*\i,1.35-0.9*\p) {\the\numexpr\i-1\relax};
}
\draw[arr] (prod.east) -- (0.8,0);
\draw[omThm,very thick] (2.3+0.62*4.5,1.7) -- (2.3+0.62*4.5,1.0) node[above,pos=0,font=\scriptsize]{validée 4};
\draw[omThm,very thick] (2.3+0.62*3.5,-1.0) -- (2.3+0.62*3.5,-1.7) node[below,pos=1,font=\scriptsize]{validée 3};
\node[draw=omThm,rounded corners=2pt,align=center,text width=24mm] (c1) at (9.8,0.9) {consommateur 1\\partitions 0, 1};
\node[draw=omThm,rounded corners=2pt,align=center,text width=24mm] (c2) at (9.8,-0.9) {consommateur 2\\partitions 2, 3};
\node[font=\scriptsize,omThm] at (9.8,-1.9) {un groupe de consommateurs};
\draw[arr] (6.95,1.35) -- (c1.west); \draw[arr] (5.7,0.45) -- (c1.west);
\draw[arr] (7.6,-0.45) -- (c2.west); \draw[arr] (6.35,-1.35) -- (c2.west);
\end{tikzpicture}

\omcaption{Un topic de quatre partitions, chacune étant une suite de messages en ajout seul repérés par leur position, et un \omterm{def:pl:messaging-and-event-driven-architecture:group}{groupe de consommateurs} de deux membres qui
se les partagent. Le groupe valide une position par partition (en rouge)~: celle à partir de laquelle un membre reprend après un redémarrage ou
un rééquilibrage.}\label{fig:pl:messaging-and-event-driven-architecture:log}
\end{omfigure}

\omcode{code/firm/eventlog/firm_eventlog.py}{53}{64}{Ajout dans un journal partitionné~: la clé choisit la partition, et la paire (identifiant, séquence) du producteur
déjà stockée est ignorée, de sorte qu'un renvoi ne crée pas de doublon.}\label{lst:pl:messaging-and-event-driven-architecture:append}

\section{Garanties de livraison et mythe de l'exactement une fois}

Un consommateur fait deux choses avec un lot~: il traite les messages, et il valide la position. Un plantage peut survenir entre les deux, et
l'ordre de ces deux opérations détermine ce que le plantage coûte.

\begin{definition}[Livraison au plus une fois et au moins une fois]\label{def:pl:messaging-and-event-driven-architecture:delivery}
Avec une \emph{livraison au plus une fois}\index{livraison au plus une fois} (at-most-once delivery), un message peut être perdu mais n'est jamais traité deux fois~: le consommateur valide
la position avant de traiter. Avec une \emph{livraison au moins une fois}\index{livraison au moins une fois} (at-least-once delivery), un message n'est jamais perdu mais peut être traité plus
d'une fois~: le consommateur valide après traitement, et un plantage entre les deux répète les messages non validés.
\end{definition}

\begin{definition}[Traitement exactement une fois]\label{def:pl:messaging-and-event-driven-architecture:exactly}
Le \emph{traitement exactement une fois}\index{traitement exactement une fois} (exactly-once processing) est la garantie que l'effet de chaque message sur l'état du consommateur se produit
une seule fois, quels que soient les plantages~: obtenue non pas du transport, qui ne peut livrer qu'au moins une fois, mais en rendant l'effet idempotent
(livre~13, chapitre~24) ou en stockant l'état et la position dans une même transaction atomique.
\end{definition}

L'article de 2011 était explicite~: «~Kafka only guarantees at-least-once delivery~», et une application qui se soucie des
doublons doit «~must add its own deduplication logic~». Des versions ultérieures ont ajouté un producteur idempotent et des transactions,
qui assurent l'exactement une fois entre topics Kafka~; pour toute autre destination, la documentation indique que l'exactement une fois «~generally
requires cooperation with such systems~». Le \omterm{def:pl:position-and-pnl-services:service}{service de positions} est une telle destination.

\omcode{code/firm/eventlog/firm_eventlog.py}{105}{126}{Un lot dans les trois modes~: validation avant traitement, validation après, ou application du lot et de la
position suivante dans une même transaction de base de données.}\label{lst:pl:messaging-and-event-driven-architecture:modes}

La journée du chapitre compte 10\,000 exécutions synthétiques pour dix comptes et vingt actions, réparties par compte dans quatre partitions et consommées par lots
de 100. Chaque jour, dix plantages surviennent en des points tirés uniformément parmi les étapes du consommateur --- chaque message traité, et l'instant entre
le traitement d'un lot et sa validation --- et cent jours sont simulés dans chaque mode. Pour chaque exécution, nous comptons le nombre de fois où elle a atteint le
\omterm{def:pl:position-and-pnl-services:service}{service de positions}.

\begin{omfigure}
\begin{tikzpicture}
\begin{axis}[width=11.5cm,height=5.4cm,ybar,bar width=14pt,ylabel={exécutions par 1\,000 plantages},symbolic x coords={at-most-once,at-least-once,atomic},
  xtick=data,xticklabels={validation d'abord\\(au plus une fois),validation après\\(au moins une fois),état et position\\atomique},
  xticklabel style={align=center},
  grid=major,grid style={gray!20},tick label style={font=\small},label style={font=\small},table/col sep=comma,ymin=0,ymax=62000,
  nodes near coords,every node near coord/.append style={font=\scriptsize,/pgf/number format/1000 sep={\,}},scaled y ticks=false,
  yticklabel style={/pgf/number format/1000 sep={\,}},
  legend style={font=\small,at={(0.5,1.03)},anchor=south,legend columns=2,draw=none,/tikz/every even column/.append style={column sep=4mm}}]
\addplot[fill=omThm!60,draw=omThm] table[x=mode,y=lost_per_1000]{figdata/platforms/24-messaging-and-event-driven-architecture/modes.csv};
\addplot[fill=omDef!40,draw=omDef] table[x=mode,y=dup_per_1000]{figdata/platforms/24-messaging-and-event-driven-architecture/modes.csv};
\legend{perdues,dupliquées}
\end{axis}
\end{tikzpicture}

\omcaption{Exécutions perdues et dupliquées pour 1\,000 plantages du consommateur, sur cent jours de dix plantages dans chaque mode. La validation d'abord perd environ
la moitié d'un lot par plantage (50\,323 pour 1\,000)~; la validation après en duplique autant (48\,652)~; état et \omterm{def:pl:messaging-and-event-driven-architecture:log}{position du consommateur} dans une même transaction, aucune. Données~:
\texttt{fig\_eventlog.py}.}\label{fig:pl:messaging-and-event-driven-architecture:modes}
\end{omfigure}

Les deux premiers modes sont des images miroir (\cref{fig:pl:messaging-and-event-driven-architecture:modes}). Un plantage au sein d'un lot de 100 tombe
en moyenne à mi-parcours, donc la validation d'abord perd environ 50 exécutions et la validation après en répète environ 50~; un plantage juste après le traitement et avant la
validation --- celle du hook --- répète les cent. Dans le troisième mode, les positions et la \omterm{def:pl:messaging-and-event-driven-architecture:log}{position du consommateur} suivante sont écrites dans une seule transaction SQLite~: un plantage
survenant avant la validation annule les deux écritures, le consommateur reprend à la position qui correspond à l'état, et aucune exécution n'est perdue ni comptée
deux fois sur 1\,000 plantages. Le même effet vient du traitement idempotent~: on stocke l'identifiant de chaque exécution avec l'état et on ignore les identifiants
déjà appliqués.

\begin{omfigure}
\begin{tikzpicture}
\begin{axis}[width=11.5cm,height=5.4cm,xlabel={jour, trié selon son erreur},ylabel={plus grande erreur de position (actions)},
  grid=major,grid style={gray!20},tick label style={font=\small},label style={font=\small},table/col sep=comma,
  xmin=1,xmax=100,ymin=-200,ymax=7000,yticklabel style={/pgf/number format/1000 sep={\,}}]
\addplot[omThm,thick] table[x=rank,y=at_most_once]{figdata/platforms/24-messaging-and-event-driven-architecture/worst.csv};
\addplot[omDef,thick] table[x=rank,y=at_least_once]{figdata/platforms/24-messaging-and-event-driven-architecture/worst.csv};
\addplot[black,thick] table[x=rank,y=atomic]{figdata/platforms/24-messaging-and-event-driven-architecture/worst.csv};
\node[omDef,font=\scriptsize,anchor=east] at (axis cs:88,6300) {validation après~: jusqu'à 6\,500};
\node[omThm,font=\scriptsize,anchor=north west] at (axis cs:55,3200) {validation d'abord~: jusqu'à 5\,100};
\node[font=\scriptsize,anchor=south west] at (axis cs:3,150) {atomique~: 0 chaque jour};
\end{axis}
\end{tikzpicture}

\omcaption{La plus grande erreur de position en fin de journée, tous comptes et actions confondus, pour chacun des cent jours, classés. Dix plantages par jour laissent une
pire erreur médiane de 3\,500 actions lorsqu'on valide d'abord et de 3\,800 lorsqu'on valide après, et jusqu'à 5\,100 et 6\,500~; le mode atomique termine
chaque jour exact. Données~: \texttt{fig\_eventlog.py}.}\label{fig:pl:messaging-and-event-driven-architecture:worst}
\end{omfigure}

\section{Patrons~: outbox, event sourcing, messages en échec}

Le problème du consommateur a son pendant côté producteur. Le service de \omterm{def:pl:trade-capture-and-booking:capture}{saisie des transactions} (trade capture) doit enregistrer une transaction dans sa base de données et publier
l'événement qui l'annonce. Deux écritures séparées ne peuvent pas être atomiques~: écrire puis publier, et un plantage entre les deux laisse une transaction dont personne n'est informé~;
publier puis écrire, et les consommateurs agissent sur une transaction qui n'existe pas.

\begin{definition}[Boîte d'envoi transactionnelle]\label{def:pl:messaging-and-event-driven-architecture:outbox}
Une \emph{boîte d'envoi transactionnelle}\index{boîte d'envoi transactionnelle} (transactional outbox) est une table de la base de données propre au service, où celui-ci écrit chaque événement au sein
de la même transaction que le changement qu'il annonce~; un relais distinct lit la table, publie les événements dans l'ordre et les marque comme envoyés, de sorte que
qu'un événement soit publié si et seulement si son changement a été validé --- au moins une fois, avec une clé qui permet au journal ou au consommateur d'écarter les
répétitions.
\end{definition}

\begin{omfigure}
\begin{tikzpicture}[font=\small,
  box/.style={draw=omDef,fill=omDef!8,rounded corners=2pt,align=center,minimum height=9mm,text width=24mm},
  arr/.style={-{Stealth[length=2mm]},thick}]
\node[box] (svc) {service de saisie\\des transactions};
\node[box,right=12mm of svc,yshift=6mm] (t) {table \texttt{trades}};
\node[box,right=12mm of svc,yshift=-6mm] (o) {table \texttt{outbox}};
\node[draw=omThm,dashed,rounded corners=3pt,fit=(t)(o),inner sep=4pt,label={[font=\scriptsize,omThm]above:une transaction}] (tx) {};
\node[box,right=18mm of tx] (r) {relais};
\node[box,right=18mm of r] (l) {journal, topic\\\texttt{trades}};
\draw[arr] (svc.east) -- (tx.west);
\draw[arr] (o.east) -- node[below,pos=0.55,font=\scriptsize]{non envoyé} (r.west |- o.east);
\draw[arr] (r) -- node[above,font=\scriptsize]{1. publier} (l);
\draw[arr,omThm] (r.south) to[bend left=35] node[below,font=\scriptsize]{2. marquer comme envoyé} (o.south east);
\end{tikzpicture}

\omcaption{La \omterm{def:pl:messaging-and-event-driven-architecture:outbox}{boîte d'envoi transactionnelle}~: la transaction et l'événement qui l'annonce sont validés dans une même transaction de base de données~; un relais publie
les événements non envoyés vers le journal, puis les marque comme envoyés. Un plantage entre 1 et 2 répète un événement sans jamais en perdre un, et la clé d'idempotence permet au
journal d'écarter la répétition.}\label{fig:pl:messaging-and-event-driven-architecture:outbox}
\end{omfigure}

\omcode{code/firm/eventlog/firm_eventlog.py}{171}{192}{La boîte d'envoi transactionnelle~: la transaction et son événement dans une même transaction de base de données, et le relais qui publie
chaque événement, avec son numéro dans la boîte d'envoi comme clé d'idempotence, avant de le marquer comme envoyé.}\label{lst:pl:messaging-and-event-driven-architecture:outbox}

\begin{table}[htbp]
\centering\small
\begin{tabular}{@{}>{\raggedright\arraybackslash}p{5.2cm}>{\centering\arraybackslash}p{2.8cm}>{\centering\arraybackslash}p{2.8cm}>{\centering\arraybackslash}p{2.8cm}@{}}
\toprule
conception & événements perdus & événements fantômes & événements dupliqués\\
\midrule
écrire la transaction, puis publier & 22 & 0 & 0\\
publier, puis écrire la transaction & 0 & 22 & 0\\
boîte d'envoi, relais sans clé d'idempotence & 0 & 0 & 22\\
boîte d'envoi, relais avec clé d'idempotence & 0 & 0 & 0\\
\bottomrule
\end{tabular}
\caption{2\,000 transactions publiées avec une probabilité de 1~\% de plantage au point dangereux de chaque conception~: entre les deux écritures, ou entre
la publication et le marquage comme envoyée d'une ligne de la boîte d'envoi. Un événement fantôme annonce une transaction que la base de données ne contient pas.}
\label{tab:pl:messaging-and-event-driven-architecture:outbox}
\end{table}

Le \Cref{tab:pl:messaging-and-event-driven-architecture:outbox} montre les trois modes de défaillance et la correction. La boîte d'envoi déplace le problème de
deux systèmes à un seul, et ce qui reste --- un relais qui a publié puis est tombé avant de marquer --- est une répétition, que le producteur idempotent ou
le consommateur l'élimine.

\begin{definition}[Event sourcing]\label{def:pl:messaging-and-event-driven-architecture:sourcing}
L'\emph{event sourcing}\index{event sourcing} stocke l'état d'un système comme la suite des événements qui l'ont changé, ajoutés et jamais édités, et
dérive chaque état courant et chaque état passé en les rejouant.
\end{definition}

Le magasin de transactions du \cref{ch:pl:trade-capture-and-booking} est en \omterm{def:pl:messaging-and-event-driven-architecture:sourcing}{event sourcing}, et le journal en est le lieu naturel~: un nouveau système aval
s'abonne au sujet et le rejoue depuis son début pour construire son propre état. Le prix à payer est la rétention --- le journal doit conserver chaque événement, ou
un instantané accompagné de la position qu'il couvre --- et une discipline sur les modifications du format des événements.

\begin{definition}[File des messages en échec]\label{def:pl:messaging-and-event-driven-architecture:dlq}
Une \emph{file des messages en échec}\index{file des messages en échec} (dead-letter queue) est un sujet vers lequel un consommateur envoie un message qu'il n'a pas réussi à traiter un nombre donné de
fois, avec la partition et la position d'où il provient, afin qu'un seul message mal formé n'arrête pas la partition qui le suit.
\end{definition}

Dans un \omterm{def:pl:messaging-and-event-driven-architecture:log}{journal partitionné}, un message qui échoue toujours est pire que dans une file~: le consommateur ne peut pas le sauter sans valider au-delà, de sorte que chaque
message qui le suit dans la partition attend. Le sujet des messages en échec permet au consommateur de poursuivre et laisse une trace à réparer et à rejouer ---
ce qui, pour les exécutions, doit être fait avant que quiconque ne se fie aux positions.

\begin{dated}{2026-09}{Les sémantiques de livraison de Kafka}\label{dat:pl:messaging-and-event-driven-architecture:kafka}
La documentation de Kafka distingue la \omterm{def:pl:messaging-and-event-driven-architecture:delivery}{livraison au plus une fois} (des messages peuvent être perdus, mais ne sont jamais redélivrés), la \omterm{def:pl:messaging-and-event-driven-architecture:delivery}{livraison au moins une fois} (ils ne sont jamais perdus, mais peuvent être
redélivrés) et la livraison exactement une fois~; un consommateur qui enregistre sa position avant le traitement obtient la première, après le traitement la deuxième. Depuis la version
0.11.0.0, son producteur propose une option de livraison idempotente, de sorte qu'un renvoi ne crée pas de doublons dans le journal, ainsi que des transactions
qui écrivent dans plusieurs partitions de façon atomique~; la livraison exactement une fois vers d'autres systèmes exige en général leur coopération.
\end{dated}

\section{Tutoriel~: une millième panne}\label{tut:pl:messaging-and-event-driven-architecture:tut}

\begin{tutorial}
\textbf{Objectif.} Consommer une journée d'exécutions soumise à des pannes selon trois modes et publier les transactions de trois façons. \textbf{État final~:}
la \cref{fig:pl:messaging-and-event-driven-architecture:modes} et le \cref{tab:pl:messaging-and-event-driven-architecture:outbox}.
\begin{tutsteps}
\item \textbf{Journal}~: \texttt{firm\_eventlog.Log}, un sujet de quatre partitions~; \texttt{pl\_eventlog.publish(fills())}.
\item \textbf{Une journée}~: \texttt{day(fs, mode, crash\_steps)} pour chaque mode, avec les mêmes étapes de panne~; compter les applications par exécution.
\item \textbf{Cent jours}~: \texttt{schedules(mode)}, exécutions perdues et dupliquées pour 1\,000 pannes et pire erreur de position.
\item \textbf{Boîte d'envoi}~: \texttt{publishing(design)} pour les quatre conceptions.
\item \textbf{Lettres mortes}~: placer une exécution mal formée dans une partition et la consommer avec \texttt{process\_with\_dlq}.
\end{tutsteps}
\textbf{Pour aller plus loin.} Remplacer le magasin atomique par un traitement idempotent (une table des \omterm{def:pl:position-and-pnl-services:execid}{identifiants d'exécution} appliqués) et vérifier qu'il donne le
même résultat~; rééquilibrer le groupe au milieu de la journée.
\end{tutorial}

\section{Construction~: le journal d'événements}\label{bld:pl:messaging-and-event-driven-architecture:eventlog}

\begin{build}
\bfield{Rôle} Transporter les événements de la société entre les services, dans l'ordre par clé, de façon durable, avec des consommateurs capables de transformer une \omterm{def:pl:messaging-and-event-driven-architecture:delivery}{livraison
au moins une fois} en effets exactement une fois.
\bfield{Interface} \texttt{Log} (\texttt{create}, \texttt{append}, \texttt{read}, \texttt{end}, \texttt{retain}, \texttt{commit},
\texttt{committed}), \texttt{assign}, \texttt{run\_batch}, \texttt{AtomicStore}, \texttt{Outbox}, \texttt{process\_with\_dlq}.
\bfield{Règles} Partitionner par clé~; ne jamais modifier un message~; chaque consommateur déclare son mode~; état et position atomiques, ou traitement idempotent,
pour tout consommateur dont l'état compte~; chaque événement publié par une boîte d'envoi~; un message qui continue d'échouer part vers les lettres mortes avec son
origine.
\bfield{Tests d'acceptation} \texttt{code/firm/eventlog/tests/}~: une clé dans une partition et dans l'ordre, segments et rétention~; un renvoi
ignoré~; l'affectation au groupe~; les trois modes sous une panne~; le renvoi du relais de la boîte d'envoi dédupliqué~; un message empoisonné vers les lettres mortes.
\bfield{Pour aller plus loin} Rééquilibrage du \omterm{def:pl:messaging-and-event-driven-architecture:group}{groupe de consommateurs} avec un numéro de génération qui écarte un membre périmé~; compaction par clé~; un serveur réseau.
\end{build}

\begin{omsources}
\item J.~Kreps, N.~Narkhede et J.~Rao, \emph{Kafka: a Distributed Messaging System for Log Processing}, NetDB 2011.
\item Documentation d'Apache Kafka, \emph{Design}~: sémantiques de livraison des messages.
\item One Quant Book~13, chapitres~12, 16 et~24 (tampons circulaires et consommateurs lents, multicast, traitement idempotent et séquenceurs).
\end{omsources}

\section{Exercices}

\begin{exercise}[$\star$]\label{exo:pl:messaging-and-event-driven-architecture:1}
Pourquoi les exécutions sont-elles réparties par compte plutôt qu'envoyées en tourniquet entre les partitions~?
\end{exercise}

\begin{exercise}[$\star$]\label{exo:pl:messaging-and-event-driven-architecture:2}
Un consommateur valide après traitement et plante après avoir appliqué 30 messages d'un lot de 100. Que se passe-t-il au redémarrage~?
\end{exercise}

\begin{exercise}[$\star$]\label{exo:pl:messaging-and-event-driven-architecture:3}
Huit partitions et trois consommateurs dans un groupe~: comment sont-elles affectées, et que se passe-t-il si un quatrième rejoint le groupe~?
\end{exercise}

\begin{exercise}[$\star\star$]\label{exo:pl:messaging-and-event-driven-architecture:4}
Pourquoi valider d'abord perd-il une cinquantaine d'exécutions par panne, et valider après en duplique-t-il à peu près autant~?
\end{exercise}

\begin{exercise}[$\star\star$]\label{exo:pl:messaging-and-event-driven-architecture:5}
Le \omterm{def:pl:position-and-pnl-services:service}{service de positions} garde son état en mémoire et écrit un instantané chaque minute. Où doit-il conserver sa position de consommateur~?
\end{exercise}

\begin{exercise}[$\star\star$]\label{exo:pl:messaging-and-event-driven-architecture:6}
Pourquoi le relais de la boîte d'envoi produit-il encore des répétitions, et qui les supprime~?
\end{exercise}

\begin{exercise}[$\star\star\star$]\label{exo:pl:messaging-and-event-driven-architecture:7}
\emph{Programmation.} Rendre idempotent le \omterm{def:pl:position-and-pnl-services:service}{service de positions} à \omterm{def:pl:messaging-and-event-driven-architecture:delivery}{livraison au moins une fois} en stockant les \omterm{def:pl:position-and-pnl-services:execid}{identifiants d'exécution} appliqués, puis exécuter \texttt{schedules} dessus.
\end{exercise}

\begin{exercise}[$\star\star\star$]\label{exo:pl:messaging-and-event-driven-architecture:8}
\emph{Trouver la faille.} «~Nous avons activé le \omterm{def:pl:messaging-and-event-driven-architecture:exactly}{traitement exactement une fois} dans le système de messagerie, donc les positions ne peuvent
plus être comptées deux fois.~»
\end{exercise}

\section{Problème~: double pendant onze minutes}

\begin{problem}[{Problème du week-end --- double pendant onze minutes}]\label{pb:pl:messaging-and-event-driven-architecture:1}
Le journal du chapitre, les exécutions, les calendriers de pannes et le service de \omterm{def:pl:trade-capture-and-booking:capture}{saisie des transactions}.

\medskip\noindent\textbf{Partie I --- Le journal.}
\begin{enumerate}
\item Que découple la \omterm{def:pl:messaging-and-event-driven-architecture:pubsub}{publication--abonnement}, et que coûte-t-elle~?
\item Quand choisir un courtier, et quand un transport sans courtier~?
\item Comment un message est-il adressé dans un \omterm{def:pl:messaging-and-event-driven-architecture:log}{journal partitionné}, et comment sa partition est-elle choisie~?
\item Que partage un \omterm{def:pl:messaging-and-event-driven-architecture:group}{groupe de consommateurs}, et qu'est-ce qu'un rééquilibrage~?
\item Que fait disparaître la rétention, et que doit faire un nouvel abonné s'il a perdu des événements dont il a besoin~?
\end{enumerate}

\medskip\noindent\textbf{Partie II --- Pannes.}
\begin{enumerate}[resume]
\item Quelles sont les deux étapes de la consommation d'un lot, et où une panne peut-elle survenir~?
\item Combien de pannes ont eu lieu dans chaque mode, et pourquoi en moins dans l'un d'eux~?
\item Combien d'exécutions ont été perdues ou dupliquées pour 1\,000 pannes dans chaque mode~?
\item Quelle est la pire erreur de position en fin de journée dans chaque mode, et la médiane~?
\item Comment le mode atomique évite-t-il à la fois les pertes et les doublons~?
\end{enumerate}

\medskip\noindent\textbf{Partie III --- Production.}
\begin{enumerate}[resume]
\item Pourquoi la transaction et son événement ne peuvent-ils pas être écrits de façon atomique sans \omterm{def:pl:messaging-and-event-driven-architecture:outbox}{boîte d'envoi transactionnelle}~?
\item Que produit chaque conception naïve en cas de pannes~?
\item Que laisse la boîte d'envoi, et comment est-elle supprimée~?
\item Qu'est-ce que l'\omterm{def:pl:messaging-and-event-driven-architecture:sourcing}{event sourcing}, et qu'exige-t-il du journal~?
\item Que protège une \omterm{def:pl:messaging-and-event-driven-architecture:dlq}{file des messages en échec}~?
\end{enumerate}

\medskip\noindent\textbf{Partie IV --- Le verdict.}
\begin{enumerate}[resume]
\item Énoncer le \emph{résultat nommé}~: le nombre d'exécutions dupliquées et perdues pour 1\,000 pannes avec la \omterm{def:pl:messaging-and-event-driven-architecture:delivery}{livraison au plus une fois}, au moins une fois et
un traitement idempotent dont la \omterm{def:pl:messaging-and-event-driven-architecture:log}{position du consommateur} est atomique, ainsi que l'erreur maximale de position que produit chacun d'eux.
\item Quel mode le \omterm{def:pl:position-and-pnl-services:service}{service de positions} doit-il utiliser, et lequel un tableau de bord de supervision~?
\item Que couvre le \omterm{def:pl:messaging-and-event-driven-architecture:exactly}{traitement exactement une fois} dans un produit de messagerie~?
\item Comment détecter l'erreur du crochet en quelques minutes~?
\item En une phrase~: d'où vient le \omterm{def:pl:messaging-and-event-driven-architecture:exactly}{traitement exactement une fois}~?
\end{enumerate}
\end{problem}

\section{Questions d'entretien}

\begin{interviewq}[$\star$ \iqroles{développeur}]\label{iq:pl:messaging-and-event-driven-architecture:1}
Expliquez la \omterm{def:pl:messaging-and-event-driven-architecture:delivery}{livraison au plus une fois}, au moins une fois et exactement une fois.
\end{interviewq}

\begin{interviewq}[$\star\star$ \iqroles{développeur}]\label{iq:pl:messaging-and-event-driven-architecture:2}
Votre consommateur applique des exécutions à une base de données. Comment vous assurez-vous qu'une panne ne compte jamais deux fois une exécution~?
\end{interviewq}

\begin{interviewq}[$\star\star$ \iqroles{développeur}]\label{iq:pl:messaging-and-event-driven-architecture:3}
Comment publiez-vous un événement chaque fois qu'une ligne est insérée dans votre base de données, sans jamais en perdre ni en inventer un~?
\end{interviewq}

\begin{interviewq}[$\star\star$ \iqroles{développeur}]\label{iq:pl:messaging-and-event-driven-architecture:4}
Comment gardez-vous dans l'ordre les messages d'un même compte lorsqu'un sujet a de nombreuses partitions et de nombreux consommateurs~?
\end{interviewq}

\begin{interviewq}[$\star\star\star$ \iqroles{développeur}]\label{iq:pl:messaging-and-event-driven-architecture:5}
Concevez le backbone d'événements qui transporte les transactions, les exécutions et les positions entre les services d'une société de trading.
\end{interviewq}
